\documentclass[11pt]{article}
\newcommand{\LabNumber}{4.2}
\newcommand{\LabTitle}{KiCad Schematic \& PCB}
\newcommand{\LabTopic}{Capturing the Lab 4-1 I/O circuit as a KiCad schematic and routing it into a real \emph{Arduino Uno shield} -- a two-layer PCB that stacks directly onto the Uno's own header pins -- optionally fabricated via PCBWay}
\newcommand{\LabDuration}{3 hours (in-lab)}
% Shared preamble for Computer Programming lab sheets.
% Include with % Shared preamble for Computer Programming lab sheets.
% Include with % Shared preamble for Computer Programming lab sheets.
% Include with \input{common.tex} after \documentclass{article}.

\usepackage[a4paper,margin=2.2cm]{geometry}
\usepackage[T1]{fontenc}
\usepackage{lmodern}
\usepackage{microtype}
\usepackage{enumitem}
\usepackage{xcolor}
\usepackage{tcolorbox}
\tcbuselibrary{breakable,skins}
\usepackage{listings}
\usepackage{fancyhdr}
\usepackage{titlesec}
\usepackage{amsmath}
\usepackage{array}
\usepackage{booktabs}
\usepackage{hyperref}

\hypersetup{
  colorlinks=true,
  linkcolor=black,
  urlcolor=blue!60!black,
  pdfborder={0 0 0}
}

% ---------- Course identity (override per file if needed) ----------
\providecommand{\CourseName}{Microprocessors}
\providecommand{\CourseCode}{010153521}
\providecommand{\Semester}{Semester 1/2026}
\providecommand{\LabDuration}{3 hours (in-lab)}

% ---------- Colors ----------
\definecolor{accent}{HTML}{1F4E79}
\definecolor{accentlight}{HTML}{D9E2EC}
\definecolor{warn}{HTML}{B45309}
\definecolor{ok}{HTML}{166534}
\definecolor{codebg}{HTML}{F6F8FA}
\definecolor{codeframe}{HTML}{D0D7DE}
\definecolor{codekw}{HTML}{0550AE}
\definecolor{codestr}{HTML}{0A7B2C}
\definecolor{codecom}{HTML}{6E7781}

% ---------- Code / assembly listings ----------
\lstdefinestyle{c}{
  language=C,
  basicstyle=\ttfamily\small,
  keywordstyle=\color{codekw}\bfseries,
  stringstyle=\color{codestr},
  commentstyle=\color{codecom}\itshape,
  backgroundcolor=\color{codebg},
  frame=single,
  rulecolor=\color{codeframe},
  framerule=0.4pt,
  showstringspaces=false,
  breaklines=true,
  columns=fullflexible,
  keepspaces=true,
  tabsize=4,
  xleftmargin=8pt,
  xrightmargin=4pt,
  aboveskip=6pt,
  belowskip=6pt,
}
\lstdefinestyle{asm}{
  basicstyle=\ttfamily\small,
  commentstyle=\color{codecom}\itshape,
  backgroundcolor=\color{codebg},
  frame=single,
  rulecolor=\color{codeframe},
  framerule=0.4pt,
  showstringspaces=false,
  breaklines=true,
  columns=fullflexible,
  keepspaces=true,
  tabsize=4,
  xleftmargin=8pt,
  xrightmargin=4pt,
  aboveskip=6pt,
  belowskip=6pt,
}
\lstset{style=asm}

% ---------- Section styling ----------
\titleformat{\section}{\Large\bfseries\color{accent}}{\thesection.}{0.5em}{}
\titleformat{\subsection}{\large\bfseries\color{accent!85!black}}{\thesubsection}{0.5em}{}
\titlespacing*{\section}{0pt}{12pt}{6pt}
\titlespacing*{\subsection}{0pt}{8pt}{4pt}

% ---------- Page headers/footers ----------
\pagestyle{fancy}
\fancyhf{}
\fancyhead[L]{\small\CourseName\ifx\CourseCode\empty\else\ (\CourseCode)\fi}
\fancyhead[C]{\small\Semester}
\fancyhead[R]{\small\textbf{Lab \LabNumber:} \LabTitle}
\fancyfoot[C]{\small\thepage}
\renewcommand{\headrulewidth}{0.4pt}

% ---------- Title block ----------
\newcommand{\labheader}{%
  \begin{tcolorbox}[
    enhanced,
    colback=accentlight,
    colframe=accent,
    boxrule=0.5pt,
    arc=2pt,
    left=10pt, right=10pt, top=8pt, bottom=8pt
  ]
    {\Large\bfseries\color{accent} Lab \LabNumber: \LabTitle}\\[2pt]
    {\small \CourseName\ifx\CourseCode\empty\else\ (\CourseCode)\fi\quad\textbar\quad \Semester}\\[1pt]
    {\small \textbf{Duration:} \LabDuration\quad\textbar\quad \textbf{Topic:} \LabTopic}
  \end{tcolorbox}
  \vspace{2pt}
}

% ---------- Custom environments ----------
\newtcolorbox{summary}[1][Quick Reference]{
  enhanced,breakable,
  colback=accentlight!40,colframe=accent!70,
  title={\bfseries #1},
  fonttitle=\color{white},
  attach boxed title to top left={xshift=6pt,yshift=-3pt},
  boxed title style={colback=accent,boxrule=0pt,arc=2pt},
  arc=2pt,boxrule=0.4pt,
  top=12pt,
}

\newcounter{labex}[section]
\newtcolorbox{labexercise}[1]{
  enhanced,breakable,
  colback=white,colframe=accent,
  title={\bfseries In-Lab Exercise \refstepcounter{labex}\thelabex: #1},
  fonttitle=\color{white},
  attach boxed title to top left={xshift=6pt,yshift=-3pt},
  boxed title style={colback=accent,boxrule=0pt,arc=2pt},
  arc=2pt,boxrule=0.4pt,
  top=12pt,
}

\newcounter{traceex}
\newtcolorbox{trace}[1][]{
  enhanced,breakable,
  colback=white,colframe=warn,
  title={\bfseries Trace \& Predict \refstepcounter{traceex}\thetraceex\ifx\\#1\\\else: #1\fi},
  fonttitle=\color{white},
  attach boxed title to top left={xshift=6pt,yshift=-3pt},
  boxed title style={colback=warn,boxrule=0pt,arc=2pt},
  arc=2pt,boxrule=0.4pt,
  top=12pt,
}

\newcounter{bugex}
\newtcolorbox{bug}[1][]{
  enhanced,breakable,
  colback=white,colframe=warn,
  title={\bfseries Find the Bug \refstepcounter{bugex}\thebugex\ifx\\#1\\\else: #1\fi},
  fonttitle=\color{white},
  attach boxed title to top left={xshift=6pt,yshift=-3pt},
  boxed title style={colback=warn,boxrule=0pt,arc=2pt},
  arc=2pt,boxrule=0.4pt,
  top=12pt,
}

\newcounter{writeex}
\newtcolorbox{writecode}[1][]{
  enhanced,breakable,
  colback=white,colframe=warn,
  title={\bfseries Write Code on Paper \refstepcounter{writeex}\thewriteex\ifx\\#1\\\else: #1\fi},
  fonttitle=\color{white},
  attach boxed title to top left={xshift=6pt,yshift=-3pt},
  boxed title style={colback=warn,boxrule=0pt,arc=2pt},
  arc=2pt,boxrule=0.4pt,
  top=12pt,
}

% ---------- AI Collaboration box ----------
\definecolor{aipurple}{HTML}{4C1D95}
\definecolor{aipurplelight}{HTML}{EDE9FE}

\newtcolorbox{aicollab}[1][AI Collaboration]{
  enhanced,breakable,
  colback=aipurplelight,colframe=aipurple!70,
  title={\bfseries #1},
  fonttitle=\color{white},
  attach boxed title to top left={xshift=6pt,yshift=-3pt},
  boxed title style={colback=aipurple,boxrule=0pt,arc=2pt},
  arc=2pt,boxrule=0.4pt,
  top=12pt,
}

\newcounter{tryex}
\newtcolorbox{tryit}[1]{
  enhanced,breakable,
  colback=ok!6,colframe=ok!70!black,
  title={\bfseries Try It \refstepcounter{tryex}\thetryex: #1},
  fonttitle=\color{white},
  attach boxed title to top left={xshift=6pt,yshift=-3pt},
  boxed title style={colback=ok!70!black,boxrule=0pt,arc=2pt},
  arc=2pt,boxrule=0.4pt,
  top=12pt,
}

% ---------- Drawing exercise ----------
\definecolor{drawteal}{HTML}{0D7377}
\definecolor{drawteallight}{HTML}{D4F1F4}

\newcounter{drawex}
\newtcolorbox{drawex}[1]{
  enhanced,breakable,
  colback=drawteallight,colframe=drawteal!70,
  title={\bfseries Draw on Paper \refstepcounter{drawex}\thedrawex: #1},
  fonttitle=\color{white},
  attach boxed title to top left={xshift=6pt,yshift=-3pt},
  boxed title style={colback=drawteal,boxrule=0pt,arc=2pt},
  arc=2pt,boxrule=0.4pt,
  top=12pt,
}


\newcommand{\takehomebanner}{%
  \vspace{6pt}
  \begin{tcolorbox}[
    colback=warn!10,colframe=warn,boxrule=0.5pt,arc=2pt,
    left=10pt,right=10pt,top=6pt,bottom=6pt
  ]
    \textbf{Take-Home (Paper Work).}\;
    Solve the following on paper without using a computer or compiler.
    Hand-write your answers; submit at the start of the next lab.
    Goal: prepare you to dry-code during the written exam.
  \end{tcolorbox}
  \vspace{4pt}
}
 after \documentclass{article}.

\usepackage[a4paper,margin=2.2cm]{geometry}
\usepackage[T1]{fontenc}
\usepackage{lmodern}
\usepackage{microtype}
\usepackage{enumitem}
\usepackage{xcolor}
\usepackage{tcolorbox}
\tcbuselibrary{breakable,skins}
\usepackage{listings}
\usepackage{fancyhdr}
\usepackage{titlesec}
\usepackage{amsmath}
\usepackage{array}
\usepackage{booktabs}
\usepackage{hyperref}

\hypersetup{
  colorlinks=true,
  linkcolor=black,
  urlcolor=blue!60!black,
  pdfborder={0 0 0}
}

% ---------- Course identity (override per file if needed) ----------
\providecommand{\CourseName}{Microprocessors}
\providecommand{\CourseCode}{010153521}
\providecommand{\Semester}{Semester 1/2026}
\providecommand{\LabDuration}{3 hours (in-lab)}

% ---------- Colors ----------
\definecolor{accent}{HTML}{1F4E79}
\definecolor{accentlight}{HTML}{D9E2EC}
\definecolor{warn}{HTML}{B45309}
\definecolor{ok}{HTML}{166534}
\definecolor{codebg}{HTML}{F6F8FA}
\definecolor{codeframe}{HTML}{D0D7DE}
\definecolor{codekw}{HTML}{0550AE}
\definecolor{codestr}{HTML}{0A7B2C}
\definecolor{codecom}{HTML}{6E7781}

% ---------- Code / assembly listings ----------
\lstdefinestyle{c}{
  language=C,
  basicstyle=\ttfamily\small,
  keywordstyle=\color{codekw}\bfseries,
  stringstyle=\color{codestr},
  commentstyle=\color{codecom}\itshape,
  backgroundcolor=\color{codebg},
  frame=single,
  rulecolor=\color{codeframe},
  framerule=0.4pt,
  showstringspaces=false,
  breaklines=true,
  columns=fullflexible,
  keepspaces=true,
  tabsize=4,
  xleftmargin=8pt,
  xrightmargin=4pt,
  aboveskip=6pt,
  belowskip=6pt,
}
\lstdefinestyle{asm}{
  basicstyle=\ttfamily\small,
  commentstyle=\color{codecom}\itshape,
  backgroundcolor=\color{codebg},
  frame=single,
  rulecolor=\color{codeframe},
  framerule=0.4pt,
  showstringspaces=false,
  breaklines=true,
  columns=fullflexible,
  keepspaces=true,
  tabsize=4,
  xleftmargin=8pt,
  xrightmargin=4pt,
  aboveskip=6pt,
  belowskip=6pt,
}
\lstset{style=asm}

% ---------- Section styling ----------
\titleformat{\section}{\Large\bfseries\color{accent}}{\thesection.}{0.5em}{}
\titleformat{\subsection}{\large\bfseries\color{accent!85!black}}{\thesubsection}{0.5em}{}
\titlespacing*{\section}{0pt}{12pt}{6pt}
\titlespacing*{\subsection}{0pt}{8pt}{4pt}

% ---------- Page headers/footers ----------
\pagestyle{fancy}
\fancyhf{}
\fancyhead[L]{\small\CourseName\ifx\CourseCode\empty\else\ (\CourseCode)\fi}
\fancyhead[C]{\small\Semester}
\fancyhead[R]{\small\textbf{Lab \LabNumber:} \LabTitle}
\fancyfoot[C]{\small\thepage}
\renewcommand{\headrulewidth}{0.4pt}

% ---------- Title block ----------
\newcommand{\labheader}{%
  \begin{tcolorbox}[
    enhanced,
    colback=accentlight,
    colframe=accent,
    boxrule=0.5pt,
    arc=2pt,
    left=10pt, right=10pt, top=8pt, bottom=8pt
  ]
    {\Large\bfseries\color{accent} Lab \LabNumber: \LabTitle}\\[2pt]
    {\small \CourseName\ifx\CourseCode\empty\else\ (\CourseCode)\fi\quad\textbar\quad \Semester}\\[1pt]
    {\small \textbf{Duration:} \LabDuration\quad\textbar\quad \textbf{Topic:} \LabTopic}
  \end{tcolorbox}
  \vspace{2pt}
}

% ---------- Custom environments ----------
\newtcolorbox{summary}[1][Quick Reference]{
  enhanced,breakable,
  colback=accentlight!40,colframe=accent!70,
  title={\bfseries #1},
  fonttitle=\color{white},
  attach boxed title to top left={xshift=6pt,yshift=-3pt},
  boxed title style={colback=accent,boxrule=0pt,arc=2pt},
  arc=2pt,boxrule=0.4pt,
  top=12pt,
}

\newcounter{labex}[section]
\newtcolorbox{labexercise}[1]{
  enhanced,breakable,
  colback=white,colframe=accent,
  title={\bfseries In-Lab Exercise \refstepcounter{labex}\thelabex: #1},
  fonttitle=\color{white},
  attach boxed title to top left={xshift=6pt,yshift=-3pt},
  boxed title style={colback=accent,boxrule=0pt,arc=2pt},
  arc=2pt,boxrule=0.4pt,
  top=12pt,
}

\newcounter{traceex}
\newtcolorbox{trace}[1][]{
  enhanced,breakable,
  colback=white,colframe=warn,
  title={\bfseries Trace \& Predict \refstepcounter{traceex}\thetraceex\ifx\\#1\\\else: #1\fi},
  fonttitle=\color{white},
  attach boxed title to top left={xshift=6pt,yshift=-3pt},
  boxed title style={colback=warn,boxrule=0pt,arc=2pt},
  arc=2pt,boxrule=0.4pt,
  top=12pt,
}

\newcounter{bugex}
\newtcolorbox{bug}[1][]{
  enhanced,breakable,
  colback=white,colframe=warn,
  title={\bfseries Find the Bug \refstepcounter{bugex}\thebugex\ifx\\#1\\\else: #1\fi},
  fonttitle=\color{white},
  attach boxed title to top left={xshift=6pt,yshift=-3pt},
  boxed title style={colback=warn,boxrule=0pt,arc=2pt},
  arc=2pt,boxrule=0.4pt,
  top=12pt,
}

\newcounter{writeex}
\newtcolorbox{writecode}[1][]{
  enhanced,breakable,
  colback=white,colframe=warn,
  title={\bfseries Write Code on Paper \refstepcounter{writeex}\thewriteex\ifx\\#1\\\else: #1\fi},
  fonttitle=\color{white},
  attach boxed title to top left={xshift=6pt,yshift=-3pt},
  boxed title style={colback=warn,boxrule=0pt,arc=2pt},
  arc=2pt,boxrule=0.4pt,
  top=12pt,
}

% ---------- AI Collaboration box ----------
\definecolor{aipurple}{HTML}{4C1D95}
\definecolor{aipurplelight}{HTML}{EDE9FE}

\newtcolorbox{aicollab}[1][AI Collaboration]{
  enhanced,breakable,
  colback=aipurplelight,colframe=aipurple!70,
  title={\bfseries #1},
  fonttitle=\color{white},
  attach boxed title to top left={xshift=6pt,yshift=-3pt},
  boxed title style={colback=aipurple,boxrule=0pt,arc=2pt},
  arc=2pt,boxrule=0.4pt,
  top=12pt,
}

\newcounter{tryex}
\newtcolorbox{tryit}[1]{
  enhanced,breakable,
  colback=ok!6,colframe=ok!70!black,
  title={\bfseries Try It \refstepcounter{tryex}\thetryex: #1},
  fonttitle=\color{white},
  attach boxed title to top left={xshift=6pt,yshift=-3pt},
  boxed title style={colback=ok!70!black,boxrule=0pt,arc=2pt},
  arc=2pt,boxrule=0.4pt,
  top=12pt,
}

% ---------- Drawing exercise ----------
\definecolor{drawteal}{HTML}{0D7377}
\definecolor{drawteallight}{HTML}{D4F1F4}

\newcounter{drawex}
\newtcolorbox{drawex}[1]{
  enhanced,breakable,
  colback=drawteallight,colframe=drawteal!70,
  title={\bfseries Draw on Paper \refstepcounter{drawex}\thedrawex: #1},
  fonttitle=\color{white},
  attach boxed title to top left={xshift=6pt,yshift=-3pt},
  boxed title style={colback=drawteal,boxrule=0pt,arc=2pt},
  arc=2pt,boxrule=0.4pt,
  top=12pt,
}


\newcommand{\takehomebanner}{%
  \vspace{6pt}
  \begin{tcolorbox}[
    colback=warn!10,colframe=warn,boxrule=0.5pt,arc=2pt,
    left=10pt,right=10pt,top=6pt,bottom=6pt
  ]
    \textbf{Take-Home (Paper Work).}\;
    Solve the following on paper without using a computer or compiler.
    Hand-write your answers; submit at the start of the next lab.
    Goal: prepare you to dry-code during the written exam.
  \end{tcolorbox}
  \vspace{4pt}
}
 after \documentclass{article}.

\usepackage[a4paper,margin=2.2cm]{geometry}
\usepackage[T1]{fontenc}
\usepackage{lmodern}
\usepackage{microtype}
\usepackage{enumitem}
\usepackage{xcolor}
\usepackage{tcolorbox}
\tcbuselibrary{breakable,skins}
\usepackage{listings}
\usepackage{fancyhdr}
\usepackage{titlesec}
\usepackage{amsmath}
\usepackage{array}
\usepackage{booktabs}
\usepackage{hyperref}

\hypersetup{
  colorlinks=true,
  linkcolor=black,
  urlcolor=blue!60!black,
  pdfborder={0 0 0}
}

% ---------- Course identity (override per file if needed) ----------
\providecommand{\CourseName}{Microprocessors}
\providecommand{\CourseCode}{010153521}
\providecommand{\Semester}{Semester 1/2026}
\providecommand{\LabDuration}{3 hours (in-lab)}

% ---------- Colors ----------
\definecolor{accent}{HTML}{1F4E79}
\definecolor{accentlight}{HTML}{D9E2EC}
\definecolor{warn}{HTML}{B45309}
\definecolor{ok}{HTML}{166534}
\definecolor{codebg}{HTML}{F6F8FA}
\definecolor{codeframe}{HTML}{D0D7DE}
\definecolor{codekw}{HTML}{0550AE}
\definecolor{codestr}{HTML}{0A7B2C}
\definecolor{codecom}{HTML}{6E7781}

% ---------- Code / assembly listings ----------
\lstdefinestyle{c}{
  language=C,
  basicstyle=\ttfamily\small,
  keywordstyle=\color{codekw}\bfseries,
  stringstyle=\color{codestr},
  commentstyle=\color{codecom}\itshape,
  backgroundcolor=\color{codebg},
  frame=single,
  rulecolor=\color{codeframe},
  framerule=0.4pt,
  showstringspaces=false,
  breaklines=true,
  columns=fullflexible,
  keepspaces=true,
  tabsize=4,
  xleftmargin=8pt,
  xrightmargin=4pt,
  aboveskip=6pt,
  belowskip=6pt,
}
\lstdefinestyle{asm}{
  basicstyle=\ttfamily\small,
  commentstyle=\color{codecom}\itshape,
  backgroundcolor=\color{codebg},
  frame=single,
  rulecolor=\color{codeframe},
  framerule=0.4pt,
  showstringspaces=false,
  breaklines=true,
  columns=fullflexible,
  keepspaces=true,
  tabsize=4,
  xleftmargin=8pt,
  xrightmargin=4pt,
  aboveskip=6pt,
  belowskip=6pt,
}
\lstset{style=asm}

% ---------- Section styling ----------
\titleformat{\section}{\Large\bfseries\color{accent}}{\thesection.}{0.5em}{}
\titleformat{\subsection}{\large\bfseries\color{accent!85!black}}{\thesubsection}{0.5em}{}
\titlespacing*{\section}{0pt}{12pt}{6pt}
\titlespacing*{\subsection}{0pt}{8pt}{4pt}

% ---------- Page headers/footers ----------
\pagestyle{fancy}
\fancyhf{}
\fancyhead[L]{\small\CourseName\ifx\CourseCode\empty\else\ (\CourseCode)\fi}
\fancyhead[C]{\small\Semester}
\fancyhead[R]{\small\textbf{Lab \LabNumber:} \LabTitle}
\fancyfoot[C]{\small\thepage}
\renewcommand{\headrulewidth}{0.4pt}

% ---------- Title block ----------
\newcommand{\labheader}{%
  \begin{tcolorbox}[
    enhanced,
    colback=accentlight,
    colframe=accent,
    boxrule=0.5pt,
    arc=2pt,
    left=10pt, right=10pt, top=8pt, bottom=8pt
  ]
    {\Large\bfseries\color{accent} Lab \LabNumber: \LabTitle}\\[2pt]
    {\small \CourseName\ifx\CourseCode\empty\else\ (\CourseCode)\fi\quad\textbar\quad \Semester}\\[1pt]
    {\small \textbf{Duration:} \LabDuration\quad\textbar\quad \textbf{Topic:} \LabTopic}
  \end{tcolorbox}
  \vspace{2pt}
}

% ---------- Custom environments ----------
\newtcolorbox{summary}[1][Quick Reference]{
  enhanced,breakable,
  colback=accentlight!40,colframe=accent!70,
  title={\bfseries #1},
  fonttitle=\color{white},
  attach boxed title to top left={xshift=6pt,yshift=-3pt},
  boxed title style={colback=accent,boxrule=0pt,arc=2pt},
  arc=2pt,boxrule=0.4pt,
  top=12pt,
}

\newcounter{labex}[section]
\newtcolorbox{labexercise}[1]{
  enhanced,breakable,
  colback=white,colframe=accent,
  title={\bfseries In-Lab Exercise \refstepcounter{labex}\thelabex: #1},
  fonttitle=\color{white},
  attach boxed title to top left={xshift=6pt,yshift=-3pt},
  boxed title style={colback=accent,boxrule=0pt,arc=2pt},
  arc=2pt,boxrule=0.4pt,
  top=12pt,
}

\newcounter{traceex}
\newtcolorbox{trace}[1][]{
  enhanced,breakable,
  colback=white,colframe=warn,
  title={\bfseries Trace \& Predict \refstepcounter{traceex}\thetraceex\ifx\\#1\\\else: #1\fi},
  fonttitle=\color{white},
  attach boxed title to top left={xshift=6pt,yshift=-3pt},
  boxed title style={colback=warn,boxrule=0pt,arc=2pt},
  arc=2pt,boxrule=0.4pt,
  top=12pt,
}

\newcounter{bugex}
\newtcolorbox{bug}[1][]{
  enhanced,breakable,
  colback=white,colframe=warn,
  title={\bfseries Find the Bug \refstepcounter{bugex}\thebugex\ifx\\#1\\\else: #1\fi},
  fonttitle=\color{white},
  attach boxed title to top left={xshift=6pt,yshift=-3pt},
  boxed title style={colback=warn,boxrule=0pt,arc=2pt},
  arc=2pt,boxrule=0.4pt,
  top=12pt,
}

\newcounter{writeex}
\newtcolorbox{writecode}[1][]{
  enhanced,breakable,
  colback=white,colframe=warn,
  title={\bfseries Write Code on Paper \refstepcounter{writeex}\thewriteex\ifx\\#1\\\else: #1\fi},
  fonttitle=\color{white},
  attach boxed title to top left={xshift=6pt,yshift=-3pt},
  boxed title style={colback=warn,boxrule=0pt,arc=2pt},
  arc=2pt,boxrule=0.4pt,
  top=12pt,
}

% ---------- AI Collaboration box ----------
\definecolor{aipurple}{HTML}{4C1D95}
\definecolor{aipurplelight}{HTML}{EDE9FE}

\newtcolorbox{aicollab}[1][AI Collaboration]{
  enhanced,breakable,
  colback=aipurplelight,colframe=aipurple!70,
  title={\bfseries #1},
  fonttitle=\color{white},
  attach boxed title to top left={xshift=6pt,yshift=-3pt},
  boxed title style={colback=aipurple,boxrule=0pt,arc=2pt},
  arc=2pt,boxrule=0.4pt,
  top=12pt,
}

\newcounter{tryex}
\newtcolorbox{tryit}[1]{
  enhanced,breakable,
  colback=ok!6,colframe=ok!70!black,
  title={\bfseries Try It \refstepcounter{tryex}\thetryex: #1},
  fonttitle=\color{white},
  attach boxed title to top left={xshift=6pt,yshift=-3pt},
  boxed title style={colback=ok!70!black,boxrule=0pt,arc=2pt},
  arc=2pt,boxrule=0.4pt,
  top=12pt,
}

% ---------- Drawing exercise ----------
\definecolor{drawteal}{HTML}{0D7377}
\definecolor{drawteallight}{HTML}{D4F1F4}

\newcounter{drawex}
\newtcolorbox{drawex}[1]{
  enhanced,breakable,
  colback=drawteallight,colframe=drawteal!70,
  title={\bfseries Draw on Paper \refstepcounter{drawex}\thedrawex: #1},
  fonttitle=\color{white},
  attach boxed title to top left={xshift=6pt,yshift=-3pt},
  boxed title style={colback=drawteal,boxrule=0pt,arc=2pt},
  arc=2pt,boxrule=0.4pt,
  top=12pt,
}


\newcommand{\takehomebanner}{%
  \vspace{6pt}
  \begin{tcolorbox}[
    colback=warn!10,colframe=warn,boxrule=0.5pt,arc=2pt,
    left=10pt,right=10pt,top=6pt,bottom=6pt
  ]
    \textbf{Take-Home (Paper Work).}\;
    Solve the following on paper without using a computer or compiler.
    Hand-write your answers; submit at the start of the next lab.
    Goal: prepare you to dry-code during the written exam.
  \end{tcolorbox}
  \vspace{4pt}
}

\usepackage{tikz}
\definecolor{pcbgreen}{HTML}{0B6E4F}
\definecolor{pcbgreendark}{HTML}{06432F}
\definecolor{silk}{HTML}{F5F5F5}
\definecolor{copper}{HTML}{C9A227}
\definecolor{pad}{HTML}{D8D8D8}
\definecolor{compblack1}{HTML}{1A1A1A}
\definecolor{compblack2}{HTML}{000000}
\definecolor{compgray1}{HTML}{9A9A9A}
\definecolor{compgray2}{HTML}{5A5A5A}

\newcommand{\drawbox}[8]{%
  % #1,#2,#3 = x0,y0,z0   #4,#5,#6 = dx,dy,dz   #7 = top color  #8 = side color
  \fill[#7] (#1,#2,#3+#6) -- (#1+#4,#2,#3+#6) -- (#1+#4,#2+#5,#3+#6) -- (#1,#2+#5,#3+#6) -- cycle;
  \fill[#8] (#1,#2,#3) -- (#1+#4,#2,#3) -- (#1+#4,#2,#3+#6) -- (#1,#2,#3+#6) -- cycle;
  \fill[#8!55!black] (#1+#4,#2,#3) -- (#1+#4,#2+#5,#3) -- (#1+#4,#2+#5,#3+#6) -- (#1+#4,#2,#3+#6) -- cycle;
  \draw[black, thin] (#1,#2,#3+#6) -- (#1+#4,#2,#3+#6) -- (#1+#4,#2+#5,#3+#6) -- (#1,#2+#5,#3+#6) -- cycle;
  \draw[black, thin] (#1,#2,#3) -- (#1+#4,#2,#3) -- (#1+#4,#2,#3+#6) -- (#1,#2,#3+#6) -- cycle;
  \draw[black, thin] (#1+#4,#2,#3) -- (#1+#4,#2+#5,#3) -- (#1+#4,#2+#5,#3+#6) -- (#1+#4,#2,#3+#6) -- cycle;
}

\begin{document}
\labheader

\section*{Objectives}
\begin{itemize}[leftmargin=*,nosep]
  \item Capture the \emph{exact} circuit you breadboarded and tested in
        \textbf{Lab 4-1} as a schematic in \textbf{KiCad}, assign footprints, and lay
        it out into a two-layer PCB that passes Electrical Rules Check (ERC) and
        Design Rules Check (DRC).
  \item Design the board in \textbf{Arduino Uno shield form factor}: the exact
        board outline and header positions the Uno itself uses, so your finished
        PCB \emph{stacks directly} onto the Uno's own pins -- no jumper wires --
        instead of connecting through a generic pin header.
  \item Generate Gerber and drill files -- the manufacturing handoff format every
        PCB fab, including PCBWay, expects.
  \item (Optional) Submit your Gerbers to \textbf{PCBWay} and order your own
        manufactured board.
\end{itemize}

\begin{summary}[Quick Reference: The Circuit You're Capturing]
This is the same 18-pin Arduino Uno R3 circuit from Lab 4-1 -- nothing new to
design, only to re-express as a schematic and land it on the Uno's real shield
headers instead of loose jumper wires.
\begin{center}
\begin{tabular}{@{}p{1.7cm}p{1.5cm}p{2.0cm}p{2.9cm}p{4.3cm}@{}}
\toprule
\textbf{Port bits} & \textbf{Uno pins} & \textbf{Direction} & \textbf{Shield header} & \textbf{Function} \\
\midrule
\texttt{PB0-PB3} & D8-D11 & Input, pull-up & Digital (8-13) & \texttt{SW1-SW4} push buttons \\
\texttt{PB4}, \texttt{PB5} & D12, D13 & Output & Digital (8-13) & \texttt{DIG1}, \texttt{DIG2} digit-select transistor bases \\
\texttt{PC0-PC5} & A0-A5 & Output & Analog In & Segments \texttt{a,b,c,d,e,f} (shared by both digits) \\
\texttt{PD2} & D2 & Output & Digital (0-7) & Segment \texttt{g} (shared by both digits) \\
\texttt{PD3-PD7} & D3-D7 & Output & Digital (0-7) & \texttt{LED1-LED5}, individually addressable \\
-- & GND & -- & Power \emph{or} Digital (8-13) & Common ground return \\
\bottomrule
\end{tabular}
\end{center}
Every signal lands on one of the Uno's \textbf{four standard shield header
strips} -- see the box below before you place a single connector symbol.
\end{summary}

\begin{summary}[Quick Reference: Arduino Uno Shield Header Layout]
A real Uno shield does not use generic pin headers wherever convenient -- it
must reproduce the Uno's own four female header strips, in their exact
positions, so the board stacks and every pin lines up:
\begin{center}
\begin{tabular}{@{}p{3.0cm}p{1.3cm}p{9.6cm}@{}}
\toprule
\textbf{Header} & \textbf{Pins} & \textbf{Signals (in order)} \\
\midrule
Digital (8-13) & 8 & \texttt{AREF, GND, D13, D12, D11, D10, D9, D8} \\
Digital (0-7) & 8 & \texttt{D7, D6, D5, D4, D3, D2, D1(TX), D0(RX)} \\
Power & 7 & \texttt{IOREF, RESET, 3V3, 5V, GND, GND, VIN} \\
Analog In & 6 & \texttt{A0, A1, A2, A3, A4, A5} \\
\bottomrule
\end{tabular}
\end{center}
The two \textbf{Digital} headers run along one long edge of the board (top, in
the usual orientation); \textbf{Power} and \textbf{Analog In} run along the
\emph{opposite} edge (bottom). We only wire the pins this circuit actually
uses (\texttt{D2-D13}, \texttt{A0-A5}, one \texttt{GND}) -- \texttt{D0}/\texttt{D1}
(USB-serial), \texttt{AREF}, and the 3V3/5V/VIN/RESET/IOREF power pins stay
unconnected on this board, exactly as they were left unused in Lab 4-1.

\textbf{The famous gap:} the horizontal spacing between the last pin of the
Digital~(0-7) header and the first pin of the Digital~(8-13) header is
\emph{not} a clean multiple of the 0.1~in grid every other pin sits on -- an
original Arduino Uno design quirk that every compatible shield must reproduce
\emph{exactly}. Do not eyeball or round this gap: copy the position from an
authoritative Arduino Uno shield footprint/outline (KiCad's bundled Arduino
connector library, if present, or Arduino's own published Uno R3 dimensional
drawing) rather than placing a plain 0.1~in header and hoping it lines up.
\end{summary}

\section*{Bill of Materials}
\begin{center}
\begin{tabular}{@{}p{1.4cm}p{6.3cm}p{7.5cm}@{}}
\toprule
\textbf{Qty} & \textbf{Part} & \textbf{Notes} \\
\midrule
4 & Tactile push button (6~mm THT) & \texttt{SW1-SW4}, no external resistor \\
5 & LED, 3~mm or 5~mm & \texttt{LED1-LED5} \\
5 & Resistor, 330~$\Omega$ & Current limiting for \texttt{LED1-LED5} \\
2 & 7-segment display, \textbf{common cathode} & \texttt{DIGIT1} (tens), \texttt{DIGIT2} (units) \\
7 & Resistor, 330~$\Omega$ & Shared segment resistors \texttt{a-g} \\
2 & NPN transistor (e.g.\ 2N3904) & Digit-select low-side switches, \texttt{Q1}/\texttt{Q2} \\
2 & Resistor, 1~k$\Omega$ & Transistor base resistors \\
4 & Female header, 2.54~mm (8, 8, 7, and 6 positions) & The Uno's standard shield
header set (Digital 8-13, Digital 0-7, Power, Analog In) -- \emph{not} generic
pin headers; these must stack directly onto the Uno's own male header pins \\
\bottomrule
\end{tabular}
\end{center}

\subsection*{KiCad Workflow Overview}
\begin{enumerate}[leftmargin=*,nosep]
  \item \textbf{Schematic capture} (\texttt{Schematic Editor}): place a symbol for
        every part in the bill of materials above, wire them exactly as you wired
        the Lab 4-1 breadboard, and label every port-bit net that connects back to
        the Uno.
  \item \textbf{Electrical Rules Check (ERC)}: catch unconnected pins, conflicting
        outputs, and missing power symbols \emph{before} laying out copper.
  \item \textbf{Footprint assignment}: attach a physical PCB footprint (THT pad
        pattern) to every symbol.
  \item \textbf{PCB layout} (\texttt{PCB Editor}): import the netlist, place
        components, reproduce the Uno's exact shield outline and header
        positions (not a board outline of your own choosing), and route every
        trace.
  \item \textbf{Design Rules Check (DRC)}: catch clearance violations, unrouted
        nets, and footprint overlaps.
  \item \textbf{Manufacturing output}: plot Gerber files (one per copper/silkscreen/
        mask layer) and a drill file -- the format every fab house, including
        PCBWay, accepts.
\end{enumerate}

\begin{summary}[Quick Reference: Symbols You Need]
\begin{center}
\begin{tabular}{@{}p{4.5cm}p{9.7cm}@{}}
\toprule
\textbf{KiCad symbol} & \textbf{Represents} \\
\midrule
\texttt{Device:R} & Every 330~$\Omega$ and 1~k$\Omega$ resistor \\
\texttt{Device:LED} & \texttt{LED1-LED5} \\
\texttt{Device:Q\_NPN\_BCE} & Digit-select transistors Q1, Q2 \\
\texttt{Switch:SW\_Push} & \texttt{SW1-SW4} \\
A 7-segment, common-cathode display symbol (search ``7SEG'' in the library
browser) & \texttt{DIGIT1}, \texttt{DIGIT2} (confirm \emph{common cathode}, not
anode, before placing) \\
An Arduino Uno shield connector footprint/symbol (search ``Arduino'' or
``Shield''; use KiCad's bundled Arduino connector library if your install has
one) & The four female header strips from the Shield Header Layout box above --
\emph{do not} substitute generic \texttt{Connector\_Generic} pin headers, since
the shield's mechanical fit depends on the real footprint, not just matching
hole spacing by eye \\
\texttt{power:GND}, \texttt{power:+5V} & Power symbols -- required by ERC on every
sheet that uses power nets \\
\bottomrule
\end{tabular}
\end{center}
This board is a \emph{peripheral shield}, not a second computer: it carries no
microcontroller, crystal, or reset circuit of its own -- the Uno you already have
supplies all of that, exactly as it did in Lab 4-1.
\end{summary}

\section*{Example Output (Illustrative Draft)}
Figures~\ref{fig:pcb2d} and~\ref{fig:pcb3d} show, roughly, the \emph{kind} of board
you are aiming for -- a rough draft sketch, \emph{not} a prescriptive layout to
copy. Your own component placement, routing, and even board outline will (and
should) look different once you complete the exercises below; every real design
decision -- where each part goes, how each trace is routed -- is yours to make.

\begin{figure}[h!]
\centering
\begin{tikzpicture}[scale=0.72, every node/.style={font=\scriptsize}]

% board outline
\fill[pcbgreen] (0,0) rectangle (12,8);
\draw[pcbgreendark, thick] (0,0) rectangle (12,8);

% mounting holes (corners)
\foreach \x/\y in {0.4/0.4, 11.6/0.4, 0.4/7.6, 11.6/7.6} {
  \fill[pad] (\x,\y) circle (0.12);
  \draw[black] (\x,\y) circle (0.12);
}

% ---- TOP edge: the two DIGITAL headers (with the famous gap) ----
\foreach \i in {0,...,7} {
  \fill[pad] (0.8+\i*0.4,7.4) circle (0.08);
}
\node[silk] at (2.2,7.05) {\ttfamily Digital (8-13)};
\foreach \i in {0,...,7} {
  \fill[pad] (4.6+\i*0.4,7.4) circle (0.08);
}
\node[silk] at (6.0,7.05) {\ttfamily Digital (0-7)};

% ---- BOTTOM edge: POWER and ANALOG IN ----
\foreach \i in {0,...,6} {
  \fill[pad] (0.8+\i*0.4,0.6) circle (0.08);
}
\node[silk] at (2.0,0.25) {\ttfamily POWER};
\foreach \i in {0,...,5} {
  \fill[pad] (4.4+\i*0.45,0.6) circle (0.08);
}
\node[silk] at (5.5,0.25) {\ttfamily ANALOG IN};

% ---- 7-segment displays ----
\foreach \x in {1.5, 5} {
  \draw[silk, thick] (\x,4.6) rectangle (\x+3,6.8);
  \node[silk] at (\x+1.5,6.7) {\scriptsize DIGIT};
  \foreach \i in {0,...,4} {
    \fill[pad] (\x+0.3+\i*0.6,4.6) circle (0.06);
  }
}

% ---- push buttons ----
\foreach \i/\name in {0/SW1,1/SW2,2/SW3,3/SW4} {
  \draw[silk, thick] (1+\i*1.6,2.2) rectangle (2+\i*1.6,3.2);
  \node[silk] at (1.5+\i*1.6,3.35) {\name};
}

% ---- transistors ----
\foreach \i/\name in {0/Q1,1/Q2} {
  \draw[silk, thick] (9.7+\i*0.9,5.0) circle (0.3);
  \node[silk] at (9.7+\i*0.9,5.5) {\name};
}

% ---- LEDs (with resistor footprints above, on the way to the top header) ----
\foreach \i in {0,...,4} {
  \draw[silk, thick] (8.6+\i*0.72,2.2) circle (0.25);
  \draw[silk] (8.6+\i*0.72-0.15,2.65) rectangle (8.6+\i*0.72+0.15,2.95);
}
\node[silk] at (9.66,1.75) {\scriptsize LED1-LED5};

% ---- resistor footprints for the shared segment bus ----
\foreach \i in {0,...,3} {
  \draw[silk] (1.35+\i*0.55,3.7) rectangle (1.65+\i*0.55,4.0);
}

% ---- illustrative copper traces (orthogonal routing; representative only) ----
% segment bus: digit1 pins -> spine -> down to Analog In header
\foreach \i in {0,...,3} {
  \draw[copper, line width=0.7pt] (1.5+\i*0.55+0.3,4.6) -- (1.5+\i*0.55+0.3,4.15);
}
\draw[copper, line width=0.9pt] (1.8,3.85) -- (6.9,3.85);
\draw[copper, line width=0.7pt] (6.9,3.85) -- (6.9,1.4) -- (6.65,0.9);

% SW1 -> Digital(8-13) header, routed left of DIGIT1
\draw[copper, line width=0.7pt] (1.4,3.2) -- (1.4,4.0) -- (0.8,4.0) -- (0.8,7.4);

% LED1 -> Digital(0-7) header, routed right of DIGIT2
\draw[copper, line width=0.7pt] (8.6,2.45) -- (8.6,2.65);
\draw[copper, line width=0.7pt] (8.6,2.95) -- (8.6,7.4) -- (7.4,7.4);

\end{tikzpicture}
\caption{Illustrative top-copper/silkscreen preview -- \emph{not to scale, not a
routing solution to copy}. The two \texttt{Digital} headers run along the top edge
(note the gap -- see the Shield Header Layout box above), \texttt{POWER} and
\texttt{ANALOG IN} along the bottom edge. Gray pads are the Uno's own shield
header pins; the thin gold lines are illustrative traces only.}
\label{fig:pcb2d}
\end{figure}

\begin{figure}[h!]
\centering
\begin{tikzpicture}[x={(0.72cm,0cm)}, y={(0.36cm,0.24cm)}, z={(0cm,0.72cm)}, font=\scriptsize]

% board slab
\drawbox{0}{0}{0}{12}{8}{0.35}{pcbgreen}{pcbgreendark}

% top digital headers (with the famous gap) -- extended past the displays'
% footprint on purpose so a sliver stays visible instead of fully hidden
\drawbox{0.4}{7.2}{0.35}{3.6}{0.4}{0.5}{compblack1}{compblack2}
\drawbox{4.5}{7.2}{0.35}{4.0}{0.4}{0.5}{compblack1}{compblack2}

% bottom power + analog-in headers
\drawbox{0.7}{0.45}{0.35}{2.6}{0.35}{0.5}{compblack1}{compblack2}
\drawbox{4.3}{0.45}{0.35}{2.5}{0.35}{0.5}{compblack1}{compblack2}

% 7-segment displays
\drawbox{1.5}{4.6}{0.35}{3}{2.2}{0.9}{compblack1}{compblack2}
\drawbox{5}{4.6}{0.35}{3}{2.2}{0.9}{compblack1}{compblack2}
\node at (3,5.7,1.3) {\textcolor{white}{\scriptsize DIGIT1}};
\node at (6.5,5.7,1.3) {\textcolor{white}{\scriptsize DIGIT2}};

% push buttons
\foreach \i in {0,1,2,3} {
  \drawbox{1+\i*1.6}{2.2}{0.35}{1}{1}{0.55}{compgray1}{compgray2}
}

% transistors
\drawbox{9.4}{4.8}{0.35}{0.6}{0.4}{0.5}{compgray1}{compgray2}
\drawbox{10.3}{4.8}{0.35}{0.6}{0.4}{0.5}{compgray1}{compgray2}

% LEDs
\foreach \i/\col in {0/red,1/red,2/green,3/yellow,4/red} {
  \drawbox{8.4+\i*0.72}{2.0}{0.35}{0.4}{0.4}{0.55}{\col}{\col!60!black}
}

\end{tikzpicture}
\caption{Illustrative 3D-rendered preview of the same draft layout -- a simplified
stand-in for what KiCad's own 3D Viewer will show once your footprints have 3D
models assigned. The two digital headers run behind the tall displays from this
angle (only their edges peek out) -- the 2D preview above is the reliable
reference for header positions. Component colors and proportions are
approximate.}
\label{fig:pcb3d}
\end{figure}

\section*{Quick Experiments (Try It)}

\begin{tryit}{One symbol, one footprint, one look}
Before wiring the whole schematic, place a single \texttt{Device:R} symbol, open its
\textbf{Footprint} field, and assign a THT footprint (e.g.\
\texttt{Resistor\_THT:R\_Axial\_DIN0207\_L6.3mm\_D2.5mm\_P10.16mm\_Horizontal}). Switch
to the PCB Editor's footprint preview and confirm you can see the two-pad footprint
before you commit to using it across all 12 resistors in your design.
\end{tryit}

\section*{In-Lab Exercises (target: 3 hours)}

\begin{labexercise}{Schematic capture \hfill {\small \itshape (50 min)}}
Create a new KiCad project. Place and wire every symbol from the Quick Reference
table above, reproducing your Lab 4-1 breadboard circuit net-for-net: buttons to
their Digital(8-13) header pins with no external resistor, LEDs through their
resistors to their Digital(0-7) header pins, the shared segment resistor bus
feeding both 7-segment symbols in parallel and continuing to the Analog In header,
and each digit-select transistor's collector/base/emitter wired as in Lab 4-1.
Annotate all references (so every part gets a unique designator like \texttt{R1},
\texttt{Q2}) and run ERC until it reports zero errors.
\end{labexercise}

\begin{labexercise}{Footprints \hfill {\small \itshape (30 min)}}
Assign a footprint to every symbol: THT axial resistor footprints for all 12
resistors, a 3~mm or 5~mm THT LED footprint (matching your real LEDs) for
\texttt{LED1-LED5}, a \texttt{TO-92} footprint for Q1/Q2, a THT tactile-switch
footprint for \texttt{SW1-SW4}, the matching 7-segment display footprint (check your
display's datasheet for pin pitch and confirm it against your physical part), and
the four Arduino shield header footprints (Digital 8-13, Digital 0-7, Power, Analog
In) -- \emph{not} generic pin headers. Re-run ERC -- footprint assignment does not
change connectivity, so it should still be clean.
\end{labexercise}

\begin{labexercise}{PCB layout and routing \hfill {\small \itshape (60 min)}}
Update the PCB from the schematic (imports the netlist as ratsnest lines). Unlike a
freestanding board, you do \emph{not} get to choose your own outline or header
positions here: reproduce the Uno's own board outline and the exact positions of
all four shield headers (roughly 68.6~mm~$\times$~53.3~mm / 2.7~in~$\times$~2.1~in
overall -- copy the real dimensions from an authoritative Arduino shield
outline/footprint rather than measuring off a photo). Place the remaining
components (displays, buttons, LEDs, transistors) wherever fits well between the
fixed headers -- displays and their digit labels easy to read from one edge,
buttons reachable from the front. Route every trace (a 0.3-0.4~mm signal trace
width is a reasonable default for this current level) until the ratsnest is empty,
then run DRC until it reports zero errors.
\begin{itemize}[nosep,leftmargin=*]
  \item \textbf{Trace/predict:} before routing, look at your ratsnest and predict
        which net will be hardest to route without crossing others -- is it the
        shared segment bus (7 nets fanning out to 14 physical pads across 2
        displays) or something else? Route it first and see if your prediction held.
  \item Confirm your Digital(0-7)/Digital(8-13) header gap matches your reference
        outline \emph{exactly} -- do not round it to the nearest 0.1~in grid
        point (see the Shield Header Layout box's ``famous gap'' warning).
\end{itemize}
\end{labexercise}

\begin{labexercise}{Manufacturing outputs \hfill {\small \itshape (25 min)}}
Plot Gerber files for every copper, mask, and silkscreen layer your board uses, plus
a drill file (Excellon format). Open the resulting Gerbers in KiCad's Gerber Viewer
(or the PCB Editor's 3D Viewer) and visually confirm: every trace you expect to see
is present, the board outline is a single closed shape, and no copper overlaps the
board edge. Zip the Gerber and drill files together -- this zip is what a fab house
expects as a single upload.
\end{labexercise}

\subsection*{Optional: Ordering from PCBWay}
Submitting your board for real fabrication is entirely optional and at your own
expense (a handful of small boards typically costs a few dollars, plus shipping) --
it is not required to complete this lab.
\begin{enumerate}[leftmargin=*,nosep]
  \item Create an account at PCBWay and start a new PCB quote.
  \item Upload the Gerber/drill zip from Exercise 4. PCBWay's online viewer will
        render your board -- compare it against your own KiCad 3D view before
        ordering; a mismatch usually means a missing or mis-plotted layer.
  \item Default specs (2 layers, 1.6~mm thickness, HASL surface finish, standard
        solder mask color) are more than adequate for this design -- there is no
        need to change anything unless you have a specific reason to.
  \item Double-check the quantity and shipping method before paying; lead time
        (fabrication plus shipping) is typically 1-2 weeks for standard service.
\end{enumerate}
If you choose not to order, submit your KiCad project files (schematic, PCB layout,
and the Gerber/drill zip) as your deliverable instead.

\begin{aicollab}
\textbf{Try these prompts with an AI tool:}
\begin{itemize}[leftmargin=*,nosep]
  \item ``What KiCad symbol and footprint should I use for a common-cathode
        7-segment display, and how do I confirm the pinout matches my physical
        part before I order a PCB?''
  \item ``What's a reasonable trace width and clearance for a hobbyist 2-layer PCB
        carrying under 100~mA per net, and why doesn't it need to be wider?''
  \item ``What's the difference between ERC and DRC in KiCad, and why can a design
        pass one but fail the other?''
  \item ``What are the exact dimensions and header positions of an Arduino Uno R3
        shield, including the gap between the two digital headers? Where does that
        information come from originally?''
\end{itemize}

\medskip\textbf{Critical check -- verify, don't just accept:}
If an AI tool gives you a footprint, symbol, or a shield dimension, do not just
place or trust it -- open the footprint preview and measure it against your
physical part's datasheet pin pitch (or against an authoritative Arduino
dimensional drawing, for the shield outline itself) before you route a single
trace. Two 7-segment displays that look identical in a library browser can use
different pin spacing, and a plausible-sounding shield dimension repeated
confidently across many sources is not the same as a verified one.

\medskip\textbf{Know the limit:}
AI tools will confidently give you a ``safe'' trace width or clearance number
without asking what copper weight (1~oz vs.\ 2~oz) or current your specific net
actually carries. Treat any such number as a starting point to verify against your
fab's own capability table (e.g.\ PCBWay's standard 1~oz/6~mil specification), not
as a final answer.
\end{aicollab}

\takehomebanner

\begin{trace}[Manufacturability at a tighter tolerance]
A different student's layout uses 0.15~mm (6~mil) trace width and 0.15~mm (6~mil)
clearance throughout, matching PCBWay's advertised \emph{standard} minimum. Another
student's layout uses 0.5~mm (20~mil) trace width and clearance throughout. Which
design is more likely to need an \emph{advanced} (costlier) fabrication tier rather
than PCBWay's standard tier, and why does clearance -- not just trace width -- matter
to that answer?
\end{trace}

\begin{bug}[Rounding the famous gap]
A student, in a hurry, placed the Digital(0-7) and Digital(8-13) headers using two
plain 8-pin, 0.1~in-pitch connectors, and simply butted them together with a clean
0.2~in gap between the last pin of one and the first pin of the other -- reasoning
that ``it's a header, it's all 0.1~in spacing.'' DRC reports zero errors, and every
schematic net is correctly connected to the correct physical pad. Explain why DRC
cannot catch this mistake, and predict exactly what happens when the student tries
to stack the fabricated shield onto a real Arduino Uno.
\end{bug}

\begin{writecode}[Hand-write a netlist]
On paper, without opening KiCad, write out the netlist for the \texttt{DIG1}
digit-select sub-circuit (Q1, its base resistor, and everything each of Q1's three
terminals connects to) in the simple text format \texttt{net\_name: pin, pin, pin}
-- one line per net. You should end up with exactly as many nets as there are
electrically distinct connection points in that sub-circuit; state how many that is
before you start counting pins.
\end{writecode}

\begin{drawex}{Sketch your board before opening KiCad}
On paper, sketch your shield's component placement -- both displays, all four
buttons, all five LEDs, and both transistors -- \emph{before} you open the PCB
Editor. Unlike Exercise 3's board outline and header positions (which are fixed by
the Uno's own dimensions, not yours to choose), component placement is still your
decision: draw the four shield headers at their real relative positions first (as
a fixed frame you must design around), then sketch where everything else goes
inside that frame. Then build your actual layout in Exercise 3 and compare: did
your paper placement survive contact with real routing constraints, or did
something need to move? What forced the change?
\end{drawex}

\section*{Reflection}

In one short paragraph, compare this lab's PCB to Lab 4-1's breadboard. They
implement the \emph{exact same} schematic -- same parts, same nets, same firmware.
What changed between them, and what stayed identical? Consider specifically: which
mistakes (a loose jumper wire, a swapped segment pin) are trivial to fix on a
breadboard but expensive or impossible to fix once a PCB is fabricated -- and what
does that imply about how much verification (ERC, DRC, comparing footprints against
real datasheets) is worth doing \emph{before} sending Gerbers to a fab house, versus
how much verification Lab 4-1's breadboard build needed before you trusted it?

\newpage
\section*{Appendix A: KiCad and PCBWay Troubleshooting}

\begin{center}
\begin{tabular}{@{}p{5.4cm}p{6.8cm}@{}}
\toprule
\textbf{Message / symptom} & \textbf{Likely cause and fix} \\
\midrule
ERC reports ``Pin not connected'' on a symbol you know is wired &
A wire crossing another wire without a junction dot is \emph{not} electrically
connected in KiCad -- add an explicit junction, or move the wire so it meets at a
clean T or corner. \\
\addlinespace
ERC reports ``conflicting power pin'' on your shield header nets &
The shield header pins need a net label or hierarchical label matching each pin
name (e.g.\ \texttt{D8}), not a floating power symbol -- power symbols are only for
true \texttt{+5V}/\texttt{GND} rails. \\
\addlinespace
Footprint preview shows the wrong pad spacing versus your physical part &
You picked a footprint by name similarity, not by measuring your part. Measure pin
pitch with calipers (or read the datasheet) and re-select a footprint that matches
exactly -- this is the single most common reason a bring-up board doesn't fit its
parts. A 3~mm LED footprint on a board populated with real 5~mm LEDs is a classic
example: ERC never catches it (footprints carry no electrical meaning to ERC),
and the part simply won't seat in undersized pads once the board is populated. \\
\addlinespace
DRC reports ``clearance violation'' between two traces that look far apart on
screen &
Zoom in -- it is almost always a via annular ring or a trace corner, not the visible
trace body. Increase zoom before assuming DRC is wrong. \\
\addlinespace
DRC reports unrouted nets after you believe everything is routed &
Check for a net with only \emph{one} pad placed (e.g.\ a symbol you forgot to update
footprints for) -- KiCad cannot route a net that only touches one physical pad. \\
\addlinespace
Shield doesn't sit flush, or some pins don't line up, once populated &
DRC and ERC only check \emph{your} design's internal consistency -- neither one
knows what a real Uno looks like. This almost always means the board outline or
header positions were measured/rounded by eye instead of copied exactly from an
authoritative Arduino Uno shield footprint or dimensional drawing (see the
``famous gap'' warning in the Shield Header Layout box). \\
\addlinespace
PCBWay's online Gerber viewer shows a blank or missing layer &
Re-check your Plot dialog's layer selection before exporting -- a common miss is
forgetting to include both silkscreen layers or the board-outline (\texttt{Edge.Cuts})
layer, which the fab needs to know your board's physical shape. \\
\bottomrule
\end{tabular}
\end{center}

\end{document}
